% GENERATED FILE. Edit canonical lessons, then run npm run generate:books.
% canonical-source-hash: fnv1a64:f23fdd2143c8ab78
% canonical-lessons: MR-C356-loksahi, MR-C356-bahumat, MR-C356-alpasankhyak, MR-C356-janmat, MR-C356-kaul-dene

\chapter{Democracy and Majorities}
\label{ch:mr-democracy-and-majorities}
\hlchaptermodality{356}

\begin{chapteropening}
\textbf{By the end of this chapter:} \emph{I can talk about democracy, majorities and public opinion.}

Complete the last lesson of chapter 356: I can talk about democracy, majorities and public opinion.
\end{chapteropening}

\section[\texorpdfstring{\mr{लोकशाही} (lokśāhī)}{लोकशाही (lokśāhī)}]{\mr{लोकशाही} (lokśāhī) — democracy}

\label{lesson:MR-C356-loksahi}



\begin{quote}
Before the new one: say the Marathi for canvassing, then the Marathi for to be elected.
\end{quote}

\subsection*{You'll want to know: \mr{लोकशाही}}
\textbf{\mr{लोकशाही}} — \emph{lokśāhī} — \textquotedblleft{}democracy\textquotedblright{}.

\begin{quote}
\mr{लोकशाहीत} \mr{जनता} \mr{सरकार} \mr{निवडते}.
\end{quote}

\begin{quote}
\emph{lokśāhīt jantā sarkār nivaḍte.}
\end{quote}

\begin{quote}
In a democracy, the people choose the government.
\end{quote}

Feminine.

\mr{लोक} + \mr{शाही}, people + rule; the opposite is \mr{हुकूमशाही}, dictatorship.

\begin{cousinweb}
\textbf{\mr{लोकशाही}} (\emph{lokśāhī}) is a Marathi compound built from \mr{लोक}, \textquotedblleft{}people\textquotedblright{}, from Sanskrit \emph{loka} (Monier-Williams), and \mr{शाही}, \textquotedblleft{}rule\textquotedblright{}, a loan from Persian \emph{shāhī}, \textquotedblleft{}kingship\textquotedblright{}.
\end{cousinweb}

\begin{grammarlens}[title={What you've built}]
One more word for giving reasons and opinions.
\end{grammarlens}

\subsection*{Your turn}
\begin{itemize}
\raggedright
  \item \emph{Say it:} \emph{lokśāhī}
  \item \emph{Say it:} \emph{lokśāhī}, once more
  \item \emph{Say it:} \emph{lokśāhīt jantā sarkār nivaḍte.} — In a democracy, the people choose the government.
  \item \emph{Recall:} say the Marathi for canvassing, then the Marathi for to be elected, then say \emph{lokśāhī} again
\end{itemize}

\begin{tcolorbox}[breakable,colback=teal!4,colframe=teal!35!black,title={Before you move on}]
What does \mr{लोकशाही} mean? (\textquotedblleft{}Democracy\textquotedblright{}.) Say it once more.
\end{tcolorbox}
\section[\texorpdfstring{\mr{बहुमत} (bahumat)}{बहुमत (bahumat)}]{\mr{बहुमत} (bahumat) — a majority}

\label{lesson:MR-C356-bahumat}



\begin{quote}
Before the new one: say the Marathi for to be elected, then the Marathi for democracy.
\end{quote}

\subsection*{You'll want to know: \mr{बहुमत}}
\textbf{\mr{बहुमत}} — \emph{bahumat} — \textquotedblleft{}a majority\textquotedblright{}.

\begin{quote}
\mr{वर्गाच्या} \mr{मतदानात} \mr{मीराला} \mr{बहुमत} \mr{मिळालं}.
\end{quote}

\begin{quote}
\emph{vargācyā matdānāt mīrālā bahumat miḷālaṁ.}
\end{quote}

\begin{quote}
Meera got a majority in the class vote.
\end{quote}

Neuter.

\mr{बहु} + \mr{मत}, many + votes.

\begin{cousinweb}
\textbf{\mr{बहुमत}} (\emph{bahumat}) is a learned compound built from Sanskrit \emph{bahu}, \textquotedblleft{}much, many\textquotedblright{}, and \emph{mata}, \textquotedblleft{}opinion\textquotedblright{} (Monier-Williams).
\end{cousinweb}

\begin{grammarlens}[title={What you've built}]
One more word for giving reasons and opinions.
\end{grammarlens}

\subsection*{Your turn}
\begin{itemize}
\raggedright
  \item \emph{Say it:} \emph{bahumat}
  \item \emph{Say it:} \emph{bahumat}, once more
  \item \emph{Say it:} \emph{vargācyā matdānāt mīrālā bahumat miḷālaṁ.} — Meera got a majority in the class vote.
  \item \emph{Recall:} say the Marathi for to be elected, then the Marathi for democracy, then say \emph{bahumat} again
\end{itemize}

\begin{tcolorbox}[breakable,colback=teal!4,colframe=teal!35!black,title={Before you move on}]
What does \mr{बहुमत} mean? (\textquotedblleft{}A majority\textquotedblright{}.) Say it once more.
\end{tcolorbox}
\section[\texorpdfstring{\mr{अल्पसंख्याक} (alpasaṅkhyāk)}{अल्पसंख्याक (alpasaṅkhyāk)}]{\mr{अल्पसंख्याक} (alpasaṅkhyāk) — a minority (group)}

\label{lesson:MR-C356-alpasankhyak}



\begin{quote}
Before the new one: say the Marathi for democracy, then the Marathi for a majority.
\end{quote}

\subsection*{You'll want to know: \mr{अल्पसंख्याक}}
\textbf{\mr{अल्पसंख्याक}} — \emph{alpasaṅkhyāk} — \textquotedblleft{}a minority (group)\textquotedblright{}.

\begin{quote}
\mr{अल्पसंख्याकांच्या} \mr{हक्कांचं} \mr{रक्षण} \mr{झालं} \mr{पाहिजे}.
\end{quote}

\begin{quote}
\emph{alpasaṅkhyākāñcyā hakkāñcaṁ rakṣaṇ jhālaṁ pāhije.}
\end{quote}

\begin{quote}
The rights of minorities must be protected.
\end{quote}

Masculine.

Also an adjective: \mr{अल्पसंख्याक} \mr{समाज}; the opposite is \mr{बहुसंख्याक}.

\begin{cousinweb}
\textbf{\mr{अल्पसंख्याक}} (\emph{alpasaṅkhyāk}) is a learned word built from Sanskrit \emph{alpa}, \textquotedblleft{}small\textquotedblright{}, and \emph{saṁkhyā}, \textquotedblleft{}number\textquotedblright{} (Monier-Williams).
\end{cousinweb}

\begin{grammarlens}[title={What you've built}]
One more word for giving reasons and opinions.
\end{grammarlens}

\subsection*{Your turn}
\begin{itemize}
\raggedright
  \item \emph{Say it:} \emph{alpasaṅkhyāk}
  \item \emph{Say it:} \emph{alpasaṅkhyāk}, once more
  \item \emph{Say it:} \emph{alpasaṅkhyākāñcyā hakkāñcaṁ rakṣaṇ jhālaṁ pāhije.} — The rights of minorities must be protected.
  \item \emph{Recall:} say the Marathi for democracy, then the Marathi for a majority, then say \emph{alpasaṅkhyāk} again
\end{itemize}

\begin{tcolorbox}[breakable,colback=teal!4,colframe=teal!35!black,title={Before you move on}]
What does \mr{अल्पसंख्याक} mean? (\textquotedblleft{}A minority (group)\textquotedblright{}.) Say it once more.
\end{tcolorbox}
\section[\texorpdfstring{\mr{जनमत} (janmat)}{जनमत (janmat)}]{\mr{जनमत} (janmat) — public opinion}

\label{lesson:MR-C356-janmat}



\begin{quote}
Before the new one: say the Marathi for a majority, then the Marathi for a minority.
\end{quote}

\subsection*{You'll want to know: \mr{जनमत}}
\textbf{\mr{जनमत}} — \emph{janmat} — \textquotedblleft{}public opinion\textquotedblright{}.

\begin{quote}
\mr{जनमत} \mr{या} \mr{निर्णयाच्या} \mr{बाजूने} \mr{होतं}.
\end{quote}

\begin{quote}
\emph{janmat yā nirṇayācyā bājūne hotaṁ.}
\end{quote}

\begin{quote}
Public opinion was in favour of this decision.
\end{quote}

Neuter.

\mr{जन} + \mr{मत}, people + opinion; neuter, so \mr{होतं}.

\begin{cousinweb}
\textbf{\mr{जनमत}} (\emph{janmat}) is a learned compound built from Sanskrit \emph{jana}, \textquotedblleft{}people\textquotedblright{}, and \emph{mata}, \textquotedblleft{}opinion\textquotedblright{} (Monier-Williams).
\end{cousinweb}

\begin{grammarlens}[title={What you've built}]
One more word for giving reasons and opinions.
\end{grammarlens}

\subsection*{Your turn}
\begin{itemize}
\raggedright
  \item \emph{Say it:} \emph{janmat}
  \item \emph{Say it:} \emph{janmat}, once more
  \item \emph{Say it:} \emph{janmat yā nirṇayācyā bājūne hotaṁ.} — Public opinion was in favour of this decision.
  \item \emph{Recall:} say the Marathi for a majority, then the Marathi for a minority, then say \emph{janmat} again
\end{itemize}

\begin{tcolorbox}[breakable,colback=teal!4,colframe=teal!35!black,title={Before you move on}]
What does \mr{जनमत} mean? (\textquotedblleft{}Public opinion\textquotedblright{}.) Say it once more.
\end{tcolorbox}
\section[\texorpdfstring{\mr{कौल} \mr{देणे} (kaul deṇe)}{कौल देणे (kaul deṇe)}]{\mr{कौल} \mr{देणे} (kaul deṇe) — to give a mandate; to give a verdict}

\label{lesson:MR-C356-kaul-dene}



\begin{quote}
Before the new one: say the Marathi for a minority, then the Marathi for public opinion.
\end{quote}

\subsection*{You'll want to know: \mr{कौल} \mr{देणे}}
\textbf{\mr{कौल} \mr{देणे}} — \emph{kaul deṇe} — \textquotedblleft{}to give a mandate; to give a verdict\textquotedblright{}.

\begin{quote}
\mr{लोकांनी} \mr{तरुण} \mr{उमेदवाराला} \mr{कौल} \mr{दिला}.
\end{quote}

\begin{quote}
\emph{lokānnī taruṇ umedvārālā kaul dilā.}
\end{quote}

\begin{quote}
The people gave their mandate to the young candidate.
\end{quote}

Past, with -\mr{ने} (he …): \mr{त्याने} \mr{कौल} \mr{दिला}, \emph{tyāne kaul dilā}.

The one chosen takes -\mr{ला}: \mr{उमेदवाराला} \mr{कौल} \mr{देणे}.

\begin{cousinweb}
\textbf{\mr{कौल} \mr{देणे}} (\emph{kaul deṇe}) joins \mr{कौल}, \textquotedblleft{}a verdict, a word given\textquotedblright{}, a loan usually traced to Arabic \emph{qaul}, \textquotedblleft{}saying, word\textquotedblright{}, through Persian, to \mr{देणे}, \textquotedblleft{}to give\textquotedblright{}.
\end{cousinweb}

\begin{grammarlens}[title={What you've built}]
One more word for giving reasons and opinions.
\end{grammarlens}

\subsection*{Your turn}
\begin{itemize}
\raggedright
  \item \emph{Say it:} \emph{kaul deṇe}
  \item \emph{Say it:} \emph{kaul deṇe}, once more
  \item \emph{Say it:} \emph{lokānnī taruṇ umedvārālā kaul dilā.} — The people gave their mandate to the young candidate.
  \item \emph{Recall:} say the Marathi for a minority, then the Marathi for public opinion, then say \emph{kaul deṇe} again
\end{itemize}

\begin{tcolorbox}[breakable,colback=teal!4,colframe=teal!35!black,title={Before you move on}]
What does \mr{कौल} \mr{देणे} mean? (\textquotedblleft{}To give a mandate; to give a verdict\textquotedblright{}.) Say it once more.
\end{tcolorbox}
